\scanpage{95}
\mspage{82}

aux conditions de la prop. 4.6, \deleted{\unsure{et}} ainsi qu'à la condition supplémentaire à (4.6.1.)
\[
4^\circ)\quad\unsure{u|v}^*(u_!(\mathrm{TC}(\unsure{I|1}))\subset\mathrm{TC}(I).\tag*{(4.8.1)}
\]
Donc on trouve une structure de cat. à cofibr. spéciale sur $M(I)$, en posant
\[
W(I)=u^{*-1}(W(I_0))\tag*{(4.8.2)}
\]
(c'est en fait \inserted{\unsure{équiv. à|égal à}} la définition habituelle de $W(I)$) et
\[
\begin{aligned}
C(I)={}&\text{clôture }\deleted{\text{\unsure{pour}}}\ u_!(C(I_0))\ \text{pour}\\ &\text{les op. de coch. de base et de}\\ &\text{composition transfinie.}
\end{aligned}\tag*{(4.8.3)}
\]

On a \inserted{de plus,} sous ces conditions :
\[
\underset{\substack{\Vert\,\mathrm{déf}\\(C(I))_*}}{\mathrm{TF}(I)} =(u^*)^{-1}\bigl(\underset{\substack{\Vert\,\mathrm{déf}\\(C_{\inserted{I_0}}\deleted{\unsure{(F)}})_*,\ \text{et}\\\text{\unsure{définit|défini} \unsure{aussi|ainsi}}\\\text{fibre par fibre,}}}{\underline{\mathrm{TF}(I_0)}}\bigr)\tag*{(4.8.8)}
\]
\margin{Ces inclusions devraient être des égalités, mais je ne sais pas \unsure{le|les} prouver sans faire des hypothèses plus fortes, du type \unsure{\S\,9|p.~9}, \uline{cat. de Quillen spéciales}.}
% slanted
\[
\left\{\begin{aligned} u_!\bigl(\overbrace{\mathrm{TC}(I_0)}^{C(I_0)\cap W(I_0)}\bigr)&\subset\mathrm{TC}(I)\\ F(I)\overset{\mathrm{déf}}{=}\bigl(\mathrm{TC}(I)\bigr)_*&\subset(u^*)^{-1}\bigl(F(I_0)\bigr).
\end{aligned}\right.\tag*{(4.8.9)}
\]

\uline{\unsure{Dém|Preuve}} On sait déjà que $M(I)$ est stable par \inserted{petites} $\varinjlim$ \deleted{$M$ l'étant,} et que \deleted{\unsure{tout|les}} \unsure{élément|l'élément} est accessible, \unsure{comme} il résulte aussitôt des propriétés idoines de $M$. Il est clair aussi que $u^*$ commute aux \inserted{petites} $\varinjlim$. Il reste à prouver (4.6\unsure{.}1. $3^\circ$), et $4^\circ$ de (4.8.1) \unsure{\deleted{q}|{}}\unsure{i.e.}
% archive sheet number: 94
